% GENERATED FILE. Edit canonical lessons, then run npm run generate:books.
% canonical-source-hash: fnv1a64:91de6cc370dd8c4c
% canonical-lessons: LA-C312-fateor, LA-C312-punio, LA-C312-supplicium, LA-C312-exsilium, LA-C312-vinculum

\chapter{Poena et supplicium}
\label{ch:la-poena-et-supplicium}
\hlchaptermodality{312}

\begin{chapteropening}
\textbf{By the end of this chapter:} \emph{I can tell what punishment followed: to admit, to punish, execution, exile and chains.}

Complete the last lesson of chapter 312: I can tell what punishment followed: to admit, to punish, execution, exile and chains.
\end{chapteropening}

\section[fateor, fatērī]{fateor, fatērī — to admit, to confess}

\label{lesson:LA-C312-fateor}



\begin{quote}
Before the new one: say the Latin for innocent, then the Latin for fault.
\end{quote}

\subsection*{You'll want to know: fateor, fatērī}
\textbf{fateor, fatērī} — \textquotedblleft{}to admit, to confess\textquotedblright{}.

\begin{quote}
Tandem fūr fassus est sē crumēnam rapuisse.
\end{quote}

\begin{quote}
\emph{At last the thief admitted that he had snatched the purse.}
\end{quote}

\emph{fateor, fatērī, fassus sum}, a deponent, with accusative and infinitive. English \emph{confess} comes from a compound of it.

\begin{cousinweb}
\textbf{fateor} is inherited from Proto-Italic \emph{*fatēōr} (EtymDB), usually traced to a root meaning \textquotedblleft{}to speak\textquotedblright{}, also seen in \emph{fārī}, \textquotedblleft{}to speak\textquotedblright{}. English \emph{confess} comes from a compound.
\end{cousinweb}

\begin{grammarlens}[title={What you've built}]
One more word for telling, explaining and describing at B1.
\end{grammarlens}

\subsection*{Your turn}
\begin{itemize}
\raggedright
  \item \emph{Say it:} \emph{fateor, fatērī}
  \item \emph{Say it:} \emph{fateor, fatērī}, once more
  \item \emph{Say it:} a sentence of your own with \emph{fateor, fatērī}
  \item \emph{Recall:} say the Latin for innocent, then the Latin for fault, then say \emph{fateor, fatērī} again
\end{itemize}

\begin{tcolorbox}[breakable,colback=teal!4,colframe=teal!35!black,title={Before you move on}]
What does fateor, fatērī mean? (\textquotedblleft{}To admit, to confess\textquotedblright{}.) Say it once more.
\end{tcolorbox}
\section[pūniō, pūnīre]{pūniō, pūnīre — to punish}

\label{lesson:LA-C312-punio}



\begin{quote}
Before the new one: say the Latin for fault, then the Latin for to admit.
\end{quote}

\subsection*{You'll want to know: pūniō, pūnīre}
\textbf{pūniō, pūnīre} — \textquotedblleft{}to punish\textquotedblright{}.

\begin{quote}
Magister discipulōs pigrōs pūniēbat.
\end{quote}

\begin{quote}
\emph{The teacher used to punish lazy pupils.}
\end{quote}

\emph{pūniō, pūnīre, pūnīvī, pūnītum}, old form \emph{poeniō}. From \emph{poena} (taught). English \emph{punish}.

\begin{cousinweb}
\textbf{pūniō} is built on \emph{poena}, \textquotedblleft{}penalty\textquotedblright{} (EtymDB), itself a loan from Greek \emph{poinē}, \textquotedblleft{}blood-money\textquotedblright{}. English \emph{punish} comes from it through French.
\end{cousinweb}

\begin{grammarlens}[title={What you've built}]
One more word for telling, explaining and describing at B1.
\end{grammarlens}

\subsection*{Your turn}
\begin{itemize}
\raggedright
  \item \emph{Say it:} \emph{pūniō, pūnīre}
  \item \emph{Say it:} \emph{pūniō, pūnīre}, once more
  \item \emph{Say it:} a sentence of your own with \emph{pūniō, pūnīre}
  \item \emph{Recall:} say the Latin for fault, then the Latin for to admit, then say \emph{pūniō, pūnīre} again
\end{itemize}

\begin{tcolorbox}[breakable,colback=teal!4,colframe=teal!35!black,title={Before you move on}]
What does pūniō, pūnīre mean? (\textquotedblleft{}To punish\textquotedblright{}.) Say it once more.
\end{tcolorbox}
\section[supplicium, suppliciī]{supplicium, suppliciī — punishment, execution; a prayer}

\label{lesson:LA-C312-supplicium}



\begin{quote}
Before the new one: say the Latin for to admit, then the Latin for to punish.
\end{quote}

\subsection*{You'll want to know: supplicium, suppliciī}
\textbf{supplicium, suppliciī} — \textquotedblleft{}punishment, execution; a prayer\textquotedblright{}.

\begin{quote}
Prōditōrēs ad supplicium ductī sunt.
\end{quote}

\begin{quote}
\emph{The traitors were led to execution.}
\end{quote}

Neuter. First a kneeling in prayer, then the kneeling of the condemned. \emph{Summum supplicium}: the death penalty.

\begin{cousinweb}
\textbf{supplicium} is built on \emph{supplex}, \textquotedblleft{}kneeling, begging\textquotedblright{} (EtymDB): a kneeling down, in prayer or to receive punishment.
\end{cousinweb}

\begin{grammarlens}[title={What you've built}]
One more word for telling, explaining and describing at B1.
\end{grammarlens}

\subsection*{Your turn}
\begin{itemize}
\raggedright
  \item \emph{Say it:} \emph{supplicium, suppliciī}
  \item \emph{Say it:} \emph{supplicium, suppliciī}, once more
  \item \emph{Say it:} a sentence of your own with \emph{supplicium, suppliciī}
  \item \emph{Recall:} say the Latin for to admit, then the Latin for to punish, then say \emph{supplicium, suppliciī} again
\end{itemize}

\begin{tcolorbox}[breakable,colback=teal!4,colframe=teal!35!black,title={Before you move on}]
What does supplicium, suppliciī mean? (\textquotedblleft{}Punishment, execution; a prayer\textquotedblright{}.) Say it once more.
\end{tcolorbox}
\section[exsilium, exsiliī]{exsilium, exsiliī — exile, banishment}

\label{lesson:LA-C312-exsilium}



\begin{quote}
Before the new one: say the Latin for to punish, then the Latin for punishment.
\end{quote}

\subsection*{You'll want to know: exsilium, exsiliī}
\textbf{exsilium, exsiliī} — \textquotedblleft{}exile, banishment\textquotedblright{}.

\begin{quote}
Cōnsul in exsilium īre coāctus est.
\end{quote}

\begin{quote}
\emph{The consul was forced to go into exile.}
\end{quote}

Neuter, also \emph{exilium}. A Roman citizen could often choose exile before his trial ended. English \emph{exile}.

\begin{cousinweb}
\textbf{exsilium} is built on \emph{exsul}, \textquotedblleft{}an exiled person\textquotedblright{} (EtymDB), which Lewis \& Short connect either with \emph{solum}, \textquotedblleft{}soil\textquotedblright{}, or with a root meaning \textquotedblleft{}to go\textquotedblright{}; the origin is not settled. English \emph{exile} comes from it through French.
\end{cousinweb}

\begin{grammarlens}[title={What you've built}]
One more word for telling, explaining and describing at B1.
\end{grammarlens}

\subsection*{Your turn}
\begin{itemize}
\raggedright
  \item \emph{Say it:} \emph{exsilium, exsiliī}
  \item \emph{Say it:} \emph{exsilium, exsiliī}, once more
  \item \emph{Say it:} a sentence of your own with \emph{exsilium, exsiliī}
  \item \emph{Recall:} say the Latin for to punish, then the Latin for punishment, then say \emph{exsilium, exsiliī} again
\end{itemize}

\begin{tcolorbox}[breakable,colback=teal!4,colframe=teal!35!black,title={Before you move on}]
What does exsilium, exsiliī mean? (\textquotedblleft{}Exile, banishment\textquotedblright{}.) Say it once more.
\end{tcolorbox}
\section[vinculum, vinculī]{vinculum, vinculī — a chain, a bond}

\label{lesson:LA-C312-vinculum}



\begin{quote}
Before the new one: say the Latin for punishment, then the Latin for exile.
\end{quote}

\subsection*{You'll want to know: vinculum, vinculī}
\textbf{vinculum, vinculī} — \textquotedblleft{}a chain, a bond\textquotedblright{}.

\begin{quote}
Latrōnēs in vinculīs tenēbantur.
\end{quote}

\begin{quote}
\emph{The bandits were kept in chains.}
\end{quote}

Neuter, also \emph{vinclum}. \emph{In vincula conicere}: to throw into prison.

\begin{cousinweb}
\textbf{vinculum} is built on \emph{vinciō}, \textquotedblleft{}to bind, to fetter\textquotedblright{}, with the tool ending \emph{-ulum} (EtymDB): that which binds.
\end{cousinweb}

\begin{grammarlens}[title={What you've built}]
One more word for telling, explaining and describing at B1.
\end{grammarlens}

\subsection*{Your turn}
\begin{itemize}
\raggedright
  \item \emph{Say it:} \emph{vinculum, vinculī}
  \item \emph{Say it:} \emph{vinculum, vinculī}, once more
  \item \emph{Say it:} a sentence of your own with \emph{vinculum, vinculī}
  \item \emph{Recall:} say the Latin for punishment, then the Latin for exile, then say \emph{vinculum, vinculī} again
\end{itemize}

\begin{tcolorbox}[breakable,colback=teal!4,colframe=teal!35!black,title={Before you move on}]
What does vinculum, vinculī mean? (\textquotedblleft{}A chain, a bond\textquotedblright{}.) Say it once more.
\end{tcolorbox}
